% Final report, IEEE conference style (two columns, Times), built with the
% standard article class because IEEEtran/TikZ are not installed here.
% Figures use the LaTeX picture environment only.
\documentclass[10pt,twocolumn,a4paper]{article}
\usepackage[top=19mm,bottom=24mm,left=16mm,right=16mm,columnsep=5mm]{geometry}
\usepackage{mathptmx}
\usepackage{amsmath,amssymb}
\usepackage{booktabs}
\usepackage{tabularx}
\usepackage{array}
\usepackage{xcolor}
\usepackage{graphicx}
\usepackage{caption}
\usepackage{cite}
\usepackage{listings}
\usepackage[hidelinks]{hyperref}

% ---------------------------------------------------------------- IEEE look
\makeatletter
\renewcommand\thesection{\Roman{section}}
\renewcommand\thesubsection{\thesection-\Alph{subsection}}
\renewcommand\@seccntformat[1]{%
  \ifnum\pdfstrcmp{#1}{subsection}=0 \Alph{subsection}.\enspace
  \else\csname the#1\endcsname.\enspace\fi}
\renewcommand\section{\@startsection{section}{1}{\z@}%
  {2.0ex plus .5ex minus .2ex}{1.0ex}{\centering\normalfont\scshape}}
\renewcommand\subsection{\@startsection{subsection}{2}{\z@}%
  {1.5ex plus .5ex minus .2ex}{0.6ex}{\normalfont\itshape}}
\makeatother
\renewcommand\thetable{\Roman{table}}
\captionsetup[figure]{name=Fig.,labelsep=period,font=footnotesize,skip=4pt}
\captionsetup[table]{name=TABLE,labelsep=newline,font={footnotesize,sc},
  justification=centering,skip=3pt}
\setlength{\parindent}{1em}
\setlength{\parskip}{0pt}
\newcolumntype{L}{>{\raggedright\arraybackslash}X}
\lstset{basicstyle=\ttfamily\scriptsize,columns=fullflexible,frame=single,
  breaklines=true,aboveskip=4pt,belowskip=4pt}
\newcommand{\fb}[3]{\framebox(#1,#2){\scriptsize\shortstack{#3}}}

\begin{document}

\twocolumn[{%
\begin{center}
{\Huge Systolic Array Matrix Multiplication Accelerator\\[0.15em]
for PicoRV32 on FPGA}\\[1.0em]
{\large Aditi Dure (EE23BTECH11016), Krishna Kaustubh (EE23BTECH11005)}\\[0.2em]
{\normalsize Computer Architecture Course Project, Final Report, November 2026}
\end{center}
\vspace{1.2em}}]

% ---------------------------------------------------------------- abstract
\noindent\textbf{\textit{Abstract}---Systolic arrays are a simple and
efficient way to accelerate matrix multiplication, the core operation of
neural-network inference. We designed a loosely-coupled INT8 matrix
multiplication accelerator around a $4\times4$ output-stationary systolic
array and integrated it with an unmodified PicoRV32 RISC-V (RV32IM) processor
on a PYNQ-Z2 (Zynq-7020) FPGA. Each processing element uses a radix-4 Booth
multiplier and a carry-save accumulator built from our own adder cells, so the
accelerator uses no DSP slices. The accelerator reads its operands from a shared
dual-port memory into banked scratchpads. It tiles matrices of up to
$64\times64$ with hardware zero padding, collects results through a four-bank
result store and writes the product back, all from a single START command.
The design passes 462,796 golden-model checks and meets timing at 50\,MHz
using 4,393 LUTs, 2,137 flip-flops, 9 block RAMs and 0 DSPs (8.3\% of the
device's LUTs). Measured on the board, it computes a $64\times64\times64$
product in 457$\times$ fewer cycles than software on the same processor, with
an estimated 430$\times$ less energy.}

\vspace{0.5em}
\noindent\textbf{\textit{Index Terms}---Systolic array, Booth multiplier,
carry-save accumulation, FPGA, RISC-V, PicoRV32, hardware accelerator.}

% ================================================================ I
\section{Introduction}

Matrix multiplication dominates the run time of machine-learning inference.
For $C=AB$ with $A\in\mathbb{Z}^{M\times K}$ and $B\in\mathbb{Z}^{K\times N}$,
each output is a dot product,
\begin{equation}
C_{i,j}=\sum_{k=0}^{K-1}A_{i,k}\,B_{k,j},
\end{equation}
so the work grows as $MNK$ multiply-accumulates (MACs). A small in-order
processor spends most of its time on loads, address arithmetic and loop control
around each multiply, so it is a poor fit for this kernel.

A systolic array (SA) is a grid of processing elements (PEs), each wired only
to its neighbours. Operands are pumped through the grid, and every PE performs
one MAC per clock on data it receives from and forwards to its neighbours
\cite{kung}. Every operand that enters the array is reused by a whole row or
column of PEs, so throughput grows with the array size while memory bandwidth
does not.

This report presents a complete SA accelerator connected to a RISC-V processor
on a single FPGA, from the adder cells up to the firmware. Our main
contributions are:
\begin{itemize}\setlength\itemsep{0pt}
\item A matrix-multiplication accelerator micro-architecture with banked
on-chip scratchpads, tiles that run back to back (\emph{drain on the fly}),
and conflict-free collection of results through a four-bank result store with
a travelling tile tag;
\item A PE built from radix-4 Booth recoding and a carry-save compressor tree
made of our own half and full adders, with no carry-propagate adder in the
accumulator loop and no DSP slices;
\item Integration with an unmodified PicoRV32 core on the same FPGA through
memory-mapped registers, with firmware that times and checks both the software
and the hardware paths;
\item An on-board evaluation of cycles, resources, timing and power, compared
against a CPU-only build of the same system.
\end{itemize}

Section~\ref{sec:related} reviews related work. Section~\ref{sec:impl}
describes the implementation, Section~\ref{sec:setup} the experimental setup
and Section~\ref{sec:results} the results. Section~\ref{sec:concl} concludes.
The appendices record the feedback from our presentation, the use of LLM
assistance, the contributions of each member and how to reproduce our results.

% ================================================================ II
\section{Related Work}\label{sec:related}

Kung introduced systolic architectures in 1982 to raise computation speed
without raising memory bandwidth \cite{kung}. Google's first TPU uses a
$256\times256$ INT8 weight-stationary array \cite{tpu}. Eyeriss classified
dataflows as no-local-reuse, weight-stationary, output-stationary and
row-stationary by what each PE holds still \cite{eyeriss}. We chose output
stationary: partial sums never move, and only 8-bit operands travel between
PEs.

On the arithmetic side, Booth recoding \cite{booth} and its radix-4 form
\cite{macsorley} halve the number of partial products. Wallace-style trees of
carry-save adders reduce those products without carry propagation
\cite{wallace}. Inayat \emph{et al.} built SA PEs on radix-8 multiplication and
reported lower power and area at some cost in delay \cite{inayat}. Our PE
combines radix-4 Booth with a carry-save accumulator, so the clock-to-clock
feedback path contains no carry chain at all.

The closest work is the loosely-coupled accelerator of Pu\c{s}ca\c{s}u
\emph{et al.} \cite{puscasu}. It is configured over APB, uses DSP-based MACs
at 500\,MHz, overlaps memory transfers with computation, and was compared
against a gem5 model of a RISC-V core, reporting 3.1--145.5$\times$ speed-ups
for $8\times8$ to $64\times64$ arrays. Integrating a RISC-V core and the
accelerator on the same FPGA was listed there as future work. We take that
step at a much smaller scale: a $4\times4$ LUT-only array next to a real RISC-V
core, measured on hardware. Table~\ref{tab:cmp} summarises the differences.

\begin{table}[t]
\caption{Comparison with the reference design}\label{tab:cmp}
\centering\footnotesize
\begin{tabularx}{\columnwidth}{lLL}
\toprule
 & Pu\c{s}ca\c{s}u \emph{et al.} \cite{puscasu} & This work \\
\midrule
Array & $8\times8$ to $64\times64$ & $4\times4$ \\
MAC & DSP macro & Radix-4 Booth + CSA, LUTs \\
Data / result & INT8 / INT16 & INT8 / INT24 (exact, $K\le256$) \\
Matrix sizes & multiples of 32 & any $1\ldots64$ (zero pad) \\
Clock & 500\,MHz (core) & 50\,MHz (92\,MHz max) \\
Host & gem5 RISC-V model & PicoRV32 on the same FPGA \\
Phases & overlapped & sequential (future work) \\
Measured & simulation & on board, end to end \\
Best speed-up & 145.5$\times$ & 457.2$\times$ \\
\bottomrule
\end{tabularx}
\end{table}

% ================================================================ III
\section{Implementation}\label{sec:impl}

The RTL is written in Verilog-2001. The accelerator, \texttt{matmul\_accel},
has one 32-bit register interface for configuration and one 32-bit port into
main memory. Fig.~\ref{fig:system} shows how it is used. The CPU and the
accelerator share one address space and one dual-port block RAM. The CPU uses
port A and the accelerator port B, so neither ever stalls the other.

\begin{figure}[t]
\centering
\setlength{\unitlength}{1mm}
\begin{picture}(84,38)(0,-2)
\put(0,22){\fb{24}{12}{PicoRV32\\RV32IM, fast mul}}
\put(58,22){\fb{26}{12}{Matmul\\accelerator}}
\put(29,-2){\fb{26}{10}{64\,KB dual-port\\block RAM}}
\put(24,31){\vector(1,0){34}}\put(26,32){\tiny MMIO regs @ 0x4000\_0000}
\put(58,25){\vector(-1,0){34}}\put(30,22.6){\tiny STATUS (poll) / irq}
\put(12,22){\vector(0,-1){19}}\put(12,3){\vector(1,0){17}}
\put(13.2,9){\tiny port A}
\put(71,22){\line(0,-1){19}}\put(71,3){\vector(-1,0){16}}
\put(55,5){\vector(1,0){11}}\put(66,5){\vector(0,1){17}}
\put(72.2,9){\tiny port B}
\put(43,12){\tiny\shortstack{A, B in\\C out}}
\end{picture}
\caption{System view. The CPU writes a few registers and polls; the
accelerator reads $A$ and $B$ and writes $C$ through its own RAM port.}
\label{fig:system}
\end{figure}

\subsection{Programming model}

A job takes four steps. (1) The CPU writes $M$, $N$, $K$ and the byte
addresses of $A$, $B$ and $C$, then writes START. (2) The accelerator fetches
$A$ and $B$. (3) It computes $C$ and copies it back to RAM as a row-major
array of 32-bit integers. (4) The CPU polls STATUS until DONE is set.

The configuration registers are staged. On START, they are checked and copied
into an active set, so the CPU can prepare the next job during a run. Illegal
jobs are rejected with an error code: a dimension outside 1--64, a misaligned
C, an address past the end of RAM, or START while busy. DONE is a level that
the CPU clears, so a polling loop cannot miss it. An interrupt
(DONE $\wedge$ IRQ\_EN) is wired to PicoRV32's \texttt{irq[3]} but is not used
by the benchmark firmware. Address decoding is complete
(Table~\ref{tab:map}): any other access reads 0 and sets a sticky bus-error
flag, and byte or half-word stores to the accelerator are dropped, so a stray
byte store cannot abort a job.

\begin{table}[t]
\caption{System memory map}\label{tab:map}
\centering\footnotesize
\begin{tabularx}{\columnwidth}{llL}
\toprule
Base & Region & Contents \\
\midrule
\texttt{0x0000\_0000} & Main RAM & 64\,KB: program, data, $A$, $B$, $C$ \\
\texttt{0x4000\_0000} & Accelerator & ID, CAPS, CTRL, STATUS, ERR, JOBS, FETCH/COMPUTE/STORE cycle counters, M, N, K, A/B/C address, IRQ\_EN \\
\texttt{0x8000\_0000} & Result port & Firmware reports (to ILA / UART) \\
other & --- & Reads 0, sets \texttt{bus\_error} \\
\bottomrule
\end{tabularx}
\end{table}

\subsection{Processing element}

Each PE holds one element of $C$ (Fig.~\ref{fig:pe}). The $A$ operand
arrives from the left and is forwarded right, and $B$ arrives from above and is
forwarded down, each through one register. A MAC is accepted only when both
valid bits are set.

The multiplier recodes the signed 8-bit $B$ in radix 4: each overlapping
3-bit window selects $0$, $\pm A$ or $\pm2A$ \cite{macsorley}. This gives four
partial products instead of eight. The partial products are zero-extended. Their
negation bits and one constant that cancels the combined sign bias form a fifth
row. Together with the two accumulator vectors, these seven rows are reduced
7$\to$5$\to$4$\to$3$\to$2 by a tree of 3:2 compressors built from our own full
and half adders.

The accumulator is kept in carry-save form as a sum vector $S$ and a carry
vector $C$, so the feedback path never resolves a carry: Vivado places it in
four logic levels. The only carry-propagate addition, $S+C$, happens when a
result leaves the PE. The width is
$2\cdot8+\lceil\log_2 256\rceil=24$ bits, which is exact for up to $K=256$
products per output with plain wrap-around arithmetic.

A one-bit \emph{drain} travels with the first operand pair of the next tile.
On drain, the PE publishes the old $S+C$ with a one-clock
\texttt{result\_valid} pulse and starts the new sum from the product arriving
on the same clock. Tiles therefore run back to back with no idle clock.

\begin{figure}[t]
\centering
\setlength{\unitlength}{1mm}
\begin{picture}(84,46)(0,2)
\put(0,43.5){\vector(1,0){4}}\put(0,44.5){\tiny $a_{in}$}
\put(4,40){\fb{12}{7}{a reg}}
\put(0,31.5){\vector(1,0){4}}\put(0,32.5){\tiny $b_{in}$}
\put(4,28){\fb{12}{7}{b reg}}
\put(16,46){\vector(1,0){5}}\put(21.5,45.4){\tiny $a_{out}$ (right)}
\put(16,29){\vector(1,0){5}}\put(21.5,28.4){\tiny $b_{out}$ (down)}
\put(16,42){\vector(1,0){30}}
\put(16,32){\vector(1,0){14}}
\put(30,30){\fb{12}{6}{Booth\\recoder}}
\put(42,33){\vector(1,0){4}}
\put(46,30){\fb{14}{16}{partial\\products\\$\times4$ +\\correction}}
\put(60,38){\vector(1,0){4}}
\put(64,30){\fb{20}{16}{3:2\\compressor\\tree\\7 $\to$ 2 rows}}
\put(69,30){\vector(0,-1){12}}
\put(79,18){\vector(0,1){12}}\put(80,23){\tiny fb}
\put(64,10){\fb{20}{8}{S, C regs\\24 b each}}
\put(64,14){\vector(-1,0){8}}
\put(40,10){\fb{16}{8}{CPA $S{+}C$\\on drain}}
\put(40,14){\vector(-1,0){12}}\put(16,13.4){\tiny result}
\put(48,4){\vector(0,1){6}}\put(49.5,4.5){\tiny drain}
\end{picture}
\caption{Processing element: operand forwarding, radix-4 Booth partial
products, carry-save accumulation and the final add on drain.}
\label{fig:pe}
\end{figure}

\subsection{Array and data feed}

Sixteen PEs form the $4\times4$ mesh \texttt{systolic\_4\_4\_array}. One input
\emph{beat} is column $k$ of a four-row group of $A$ and row $k$ of a
four-column group of $B$: two 32-bit words. This is exactly what one array
edge consumes per clock, and it supports any $K$ from 1 to 256. Internal skew
registers delay row $r$ and column $c$ by $r$ and $c$ clocks so that $A_{i,k}$
and $B_{k,j}$ meet in PE$(i,j)$. A delay chain gives each PE$(r,c)$ its drain
$r+c$ clocks after PE$(0,0)$, so one tile's 16 results leave over seven
consecutive clocks.

\subsection{Memory organisation}

The \emph{fetch unit} streams $A$ and $B$ from main RAM at one word per clock
into an 8-byte buffer. It realigns rows that start at any byte address and
writes zeros to round $M$ and $N$ up to multiples of four ($K$ is never
padded). A look-ahead rule requests a word only when the buffer will have room.

Scratchpad~A has four banks, with row $r$ stored in bank $r \bmod 4$ at word
$16\lfloor r/4\rfloor+\lfloor k/4\rfloor$. The four rows of a group therefore
sit in different banks, and one address returns $A_{4g..4g+3,\,k}$ already
packed for the array. With the fixed 16-word row stride, every address is just
wires and needs no multiplier. Scratchpad~B is one RAM addressed
$\{k,\text{group}\}$, because a row slice of $B$ is four adjacent bytes. Both
scratchpads have the same one-clock read latency as main RAM. Their purpose is
bandwidth and reuse, not latency (Section~\ref{sec:analysis}).

\subsection{Tiling and result collection}

The \emph{tile sequencer} visits the output tiles with $A$ row groups in the
outer loop, $B$ column groups in the inner loop and $k$ innermost. $A$ is
re-read from its scratchpad for every column group, which is free because the
port is idle otherwise. One slot counter provides both the $k$ address and the
tile boundary, and the drain goes out with slot 0 of each new tile.

Results are written into the \emph{C store}: four banks, with row $r$ of $C$
in bank $r \bmod 4$, and up to four writes per clock. Instead of storing a tile
number in every PE, a seven-stage tag shift register carries $\{g_i,g_j\}$
beside the drain wave, and PE$(r,c)$ uses stage $r+c$. The write address
$\{g_i,g_j,c\}$ is a pure bit concatenation. For $K<4$, two PEs of one row
would publish on the same clock into the same bank. The sequencer prevents this
by stretching the tile period to $P=\max(K,4)$, inserting $4-K$ empty beats; the
array itself never stops. We found this hazard with our cycle-by-cycle
visualizer before writing the RTL. Finally, the \emph{C copier} walks the
C store row by row and writes ordinary row-major words to main RAM.

\subsection{Job control}

A job-controller FSM (IDLE, FETCH\_A, FETCH\_B, COMPUTE, COPY)
runs the job and hands RAM port~B to the fetch unit or to the copier. ABORT
works in every phase. Separate counters record fetch, compute and store cycles.
Fig.~\ref{fig:arch} shows the complete micro-architecture. Fig.~\ref{fig:phase}
shows that the phases run strictly one after another, unlike the overlapped
flow of \cite{puscasu}.

\begin{figure*}[t]
\centering
\setlength{\unitlength}{1mm}
\begin{picture}(178,64)(0,-7)
% control row
\put(0,46){\fb{24}{10}{PicoRV32\\(MMIO)}}
\put(24,51){\vector(1,0){12}}
\put(36,46){\fb{30}{10}{accel\_regs\\staging + active set, checks}}
\put(66,51){\vector(1,0){12}}
\put(78,46){\fb{64}{10}{job\_controller FSM\\FETCH\_A $\to$ FETCH\_B $\to$ COMPUTE $\to$ COPY}}
\put(110,46){\multiput(0,0)(0,-2){4}{\line(0,-1){1}}}
\multiput(41,38)(2,0){65}{\line(1,0){1}}
\put(41,38){\multiput(0,0)(0,-2){6}{\line(0,-1){1}}}
\put(114,38){\multiput(0,0)(0,-2){3}{\line(0,-1){1}}}
\put(170,38){\multiput(0,0)(0,-2){7}{\line(0,-1){1}}}
\put(128,39){\tiny phase control}
% datapath
\put(0,4){\fb{22}{22}{main RAM\\port B\\(64 KB)}}
\put(22,15){\vector(1,0){8}}
\put(30,4){\fb{22}{22}{fetch\_unit\\align, zero pad\\1 word/clock}}
\put(52,23){\vector(1,0){8}}
\put(52,7){\vector(1,0){8}}
\put(60,18){\fb{24}{10}{scratchpad\_a\\4 banks, row mod 4}}
\put(60,2){\fb{24}{10}{scratchpad\_b\\1 RAM}}
\put(92,0){\framebox(46,32){}}\put(94,29){\tiny tile\_sequencer}
\put(84,23){\vector(1,0){11}}
\put(84,7){\vector(1,0){11}}
\put(95,5){\fb{10}{20}{input\\skew\\0--3}}
\put(105,15){\vector(1,0){4}}
\put(109,5){\fb{18}{20}{$4\times4$\\PE array\\(16 PEs)}}
\put(127,15){\vector(1,0){4}}
\put(131,5){\fb{6}{20}{\rotatebox{90}{tag $\times7$}}}
\put(138,15){\vector(1,0){6}}
\put(144,8){\fb{16}{16}{c\_store\\4 banks\\$\le$4 wr/clk}}
\put(160,15){\vector(1,0){4}}
\put(164,8){\fb{14}{16}{c\_copier\\1 word\\/clock}}
\put(171,8){\line(0,-1){14}}
\put(171,-6){\line(-1,0){160}}
\put(11,-6){\vector(0,1){10}}
\put(70,-5){\tiny $C$ written back row-major (port B)}
\end{picture}
\caption{Accelerator micro-architecture. Solid arrows are data; dashed lines
are phase control from the job controller.}
\label{fig:arch}
\end{figure*}

\begin{figure}[t]
\centering
\setlength{\unitlength}{1mm}
\begin{picture}(84,30)(0,-6)
\put(0,19){\scriptsize Fetch A, B}
\put(0,12){\scriptsize Compute}
\put(0,5){\scriptsize Copy C}
\put(20,18.5){\framebox(13.0,4){\tiny 140}}
\put(33.0,11.5){\framebox(24.9,4){\tiny 268}}
\put(57.9,4.5){\framebox(24.1,4){\tiny 259}}
\put(20,1){\line(1,0){64}}
\multiput(20,1)(18.6,0){4}{\line(0,-1){1}}
\put(19,-3){\tiny 0}\put(36,-3){\tiny 200}\put(54.6,-3){\tiny 400}
\put(73.2,-3){\tiny 600}\put(70,-6){\tiny clock cycles}
\end{picture}
\caption{Phase timeline of one $16\times16\times16$ job (simulation, 667
accelerator cycles). The three phases do not overlap.}
\label{fig:phase}
\end{figure}

\begin{table*}[t]
\caption{Design alternatives considered at each layer}\label{tab:design}
\centering\footnotesize
\begin{tabularx}{\textwidth}{lp{4.3cm}p{5.4cm}L}
\toprule
Layer & Simple alternative & Chosen & Gain / trade-off \\
\midrule
Multiplier & Array / Baugh--Wooley, 8 partial products & Radix-4 Booth, 4 partial products & Half the rows, 0 DSPs; recoder and sign-correction logic \\
Accumulator & Ripple-carry add each clock & Carry-save $S,C$; one add on drain & No carry chain in loop; twice the accumulator flip-flops \\
Width & 32 bits & 24 bits $=16+\log_2256$ & Exact for $K\le256$, smaller PEs \\
Dataflow & Weight-stationary & Output-stationary & Only INT8 values move; results must be collected from 16 PEs \\
Tile boundary & Stop, read out, restart & Drain on the fly & No idle clock between tiles; results leave on 7 different clocks \\
Operands & Read main RAM directly & Banked BRAM scratchpads & One read per beat, $16\times$ reuse at $64^3$; extra fetch phase \\
Odd sizes & Mask edge PEs & Zero padding while fetching & No masking logic; up to 3 wasted rows/columns \\
Collecting $C$ & Serial scan, period $\ge16$ & 4 banks, period $\ge4$ & Up to $4\times$ faster for small $K$, same BRAM \\
Result address & Tag in every PE & 7-stage tag shift beside drain & 7 registers, no adder; needs period $\ge4$ \\
Coupling & PCPI custom instruction & Memory-mapped registers & Unmodified CPU and compiler; a few cycles of register traffic \\
Completion & Interrupt handler & Polling (IRQ built) & Exact cycle counts; CPU idles while waiting \\
\bottomrule
\end{tabularx}
\end{table*}

Table~\ref{tab:design} lists the alternatives we compared at each layer. We
evaluated the four ways of collecting $C$ cycle by cycle in our visualizer
before choosing the four-bank store.

% ================================================================ IV
\section{Experimental Setup}\label{sec:setup}

\subsection{FPGA platform}

Table~\ref{tab:platform} summarises the platform. Only the Zynq's
programmable logic is used; the ARM cores and the DDR memory are not. All
memories are on-chip block RAM inferred from one \texttt{dual\_port\_ram}
module. We built two bitstreams that are identical except for the accelerator:
the \emph{accelerated} build and a \emph{CPU-only} build
(\texttt{soc\_top ENABLE\_ACCEL=0}). Both contain the same debug cores: an ILA
to capture the result port, a VIO to configure jobs, and a UART bridge.

\begin{table}[t]
\caption{Experimental platform}\label{tab:platform}
\centering\footnotesize
\begin{tabularx}{\columnwidth}{lL}
\toprule
Board / device & TUL PYNQ-Z2, Zynq-7020 \texttt{xc7z020clg400-1} \\
Clock & 50\,MHz (MMCM from 125\,MHz) \\
Tools & Vivado 2025.2; Icarus Verilog; RISC-V GCC \\
CPU & PicoRV32, RV32IM, \texttt{ENABLE\_FAST\_MUL}, unmodified \\
Memory & 64\,KB true dual-port BRAM, 1-clock read \\
Firmware & bare-metal C, \texttt{-march=rv32im -O2} \\
Data & random INT8, seed 1; every $C$ element checked \\
\bottomrule
\end{tabularx}
\end{table}

\subsection{Baseline and metrics}

The baseline is a triple-loop matrix multiplication in C on the same
PicoRV32, using the same memory and the hardware multiplier. Both paths are
timed by the CPU with \texttt{rdcycle}, and the software loop also with
\texttt{rdinstret}. The accelerator time is \emph{end to end}. It covers the
whole driver call: register writes, START, fetch, compute, copy-back and the
polling loop. This is stricter than the metric of \cite{puscasu}, which counts
from the first scratchpad read to the last result. The VIO sets each job, the
firmware runs the accelerator and then the software, compares every element of
$C$, and reports through the result port. A Tcl script (\texttt{run.tcl})
captures the report with the ILA and decodes it into \texttt{results.csv}.

\subsection{Verification environment}

Fig.~\ref{fig:verif} shows the verification flow. Expected values come only
from a Python golden model of exact INT8 arithmetic with 24-bit wrap-around.
Every block has a self-checking testbench that prints the expected and actual
value of each check and fails on any mismatch, unknown value, missing or extra
result, or timeout (Table~\ref{tab:verif}). The board firmware is also run in
Icarus, and its result stream is decoded by the same script used on the board.

\begin{figure}[t]
\centering
\setlength{\unitlength}{1mm}
\begin{picture}(84,34)(0,0)
\put(0,22){\fb{18}{10}{golden\\model.py}}
\put(18,27){\vector(1,0){5}}
\put(23,22){\fb{16}{10}{vectors\\(A, B, C)}}
\put(39,27){\vector(1,0){5}}
\put(44,22){\fb{20}{10}{self-checking\\testbench}}
\put(64,28){\vector(1,0){4}}\put(68,26){\vector(-1,0){4}}
\put(68,22){\fb{16}{10}{DUT\\(RTL)}}
\put(54,22){\vector(0,-1){4}}
\put(44,12){\fb{20}{6}{PASS / FAIL}}
\put(0,0){\fb{24}{10}{firmware: SW\\vs accelerator}}
\put(24,5){\vector(1,0){5}}
\put(29,0){\fb{20}{10}{soc\_top or\\board (ILA)}}
\put(49,5){\vector(1,0){5}}
\put(54,0){\fb{30}{10}{check\_fpga\_run.py\\decode + check}}
\end{picture}
\caption{Verification flow: unit and accelerator testbenches (top) and
firmware-level checking in simulation and on the board (bottom).}
\label{fig:verif}
\end{figure}

\begin{table}[t]
\caption{Verification coverage}\label{tab:verif}
\centering\footnotesize
\begin{tabularx}{\columnwidth}{lLr}
\toprule
Level & What is checked & Checks \\
\midrule
Adders, Booth, tree & exhaustive small widths (all 65,536 INT8 products), random large & 89,741 \\
PE & accumulate, drain, clear, stalls & 20,980 \\
Memories, fetch & every word, every alignment & 33,449 \\
Tile sequencer & every $C$ write, exact cycles & 11,042 \\
C copier & stalls, reset mid-copy & 45,136 \\
Accelerator & golden jobs, all error codes, ABORT, IRQ & 262,333 \\
CPU system & firmware vs software & 15 jobs \\
\midrule
\multicolumn{2}{l}{\texttt{make check}: 40 testbenches, 0 failures} & \textbf{462,796} \\
\bottomrule
\end{tabularx}
\end{table}

% ================================================================ V
\section{Experimental Results}\label{sec:results}

\subsection{FPGA utilization}

Table~\ref{tab:util} compares the two routed builds. The accelerator adds
4,401 LUTs, 2,140 flip-flops and 9 BRAM tiles, and no DSPs. The four DSPs in
both builds belong to the CPU's fast multiplier. Table~\ref{tab:blocks} breaks
the cost down by block. The tile sequencer, which contains the 16 PEs, takes
76\% of the accelerator's LUTs; a single PE synthesised out of context uses 195
LUTs and 88 flip-flops. Scratchpad~A uses four RAMB36 although each bank holds
only 1\,KB, because a true dual-port RAM wider than 18 bits needs a RAMB36 on
7-series parts. The C store uses another four, and scratchpad~B one. With the
debug cores excluded, the accelerator is 3.6$\times$ the size of the CPU, but
it uses only 8.3\% of the device's LUTs, 2.0\% of its flip-flops and 6.4\% of
its BRAM.

\begin{table}[t]
\caption{Routed utilization, timing and power}\label{tab:util}
\centering\footnotesize
\begin{tabular}{lrrr}
\toprule
xc7z020 & CPU only & CPU + accel. & $\Delta$ \\
\midrule
LUTs & 2,776 (5.2\%) & 7,177 (13.5\%) & +4,401 \\
Flip-flops & 4,090 (3.8\%) & 6,230 (5.9\%) & +2,140 \\
BRAM tiles & 20 (14.3\%) & 29 (20.7\%) & +9 \\
DSPs & 4 & 4 & 0 \\
WNS @ 50\,MHz & +11.032\,ns & +9.130\,ns & \\
$f_{max}$ & $\approx$111.5\,MHz & $\approx$92.0\,MHz & $-17\%$ \\
Power (est.) & 0.228\,W & 0.242\,W & +14\,mW \\
\bottomrule
\end{tabular}
\end{table}

\begin{table}[t]
\caption{Utilization by block (accelerated build)}\label{tab:blocks}
\centering\footnotesize
\begin{tabular}{lrrrr}
\toprule
Block & LUTs & FFs & BRAM & DSP \\
\midrule
PicoRV32 (incl.\ fast mul) & 1,217 & 638 & 0 & 4 \\
\quad of which \texttt{pcpi\_fast\_mul} & 277 & 37 & 0 & 4 \\
Main RAM (64\,KB) & --- & --- & 16 & --- \\
\midrule
\texttt{tile\_sequencer} (16 PEs) & 3,353 & 1,526 & 0 & 0 \\
\texttt{fetch\_unit} & 401 & 153 & 0 & 0 \\
\texttt{accel\_regs} & 379 & 336 & 0 & 0 \\
\texttt{job\_controller} & 74 & 73 & 0 & 0 \\
\texttt{scratchpad\_b} & 80 & 0 & 1 & 0 \\
scratchpad A, C store, copier & 106 & 49 & 8 & 0 \\
\textbf{\texttt{matmul\_accel} total} & \textbf{4,393} & \textbf{2,137} & \textbf{9} & \textbf{0} \\
\midrule
Debug (ILA, VIO, hub, UART) & $\approx$1,550 & $\approx$3,300 & 4 & 0 \\
\bottomrule
\end{tabular}
\end{table}

\subsection{Timing and power}

Both builds meet timing at 50\,MHz. The accelerator lowers the estimated
$f_{max}$ from about 111.5 to about 92\,MHz. The limiting path is not the Booth
MAC. It is the drain control bit (\texttt{drain\_shift\_reg[2]}), which fans
out to 655 inputs across the 16 PEs; the path is 10.8\,ns, 92\% of it routing
delay. Out of context, a PE's carry-save loop closes at about 195\,MHz. In the
CPU-only build, the limiting path is the CPU's memory-read path (89\% routing).
Vivado estimates that the accelerator adds 14\,mW (+6\%).

\subsection{Performance}

Table~\ref{tab:perf} and Figs.~\ref{fig:cycles}--\ref{fig:speedup} give the
on-board results for square matrices. Every job passed. The CPU-only bitstream
reproduces exactly the same software cycle counts, so adding the accelerator
does not slow the processor. Software time grows by $\approx8\times$ per
doubling of $n$. The accelerator's time grows more slowly while fixed costs
dominate and approaches the same factor at large sizes (6.1$\times$ from
$32^3$ to $64^3$). The speed-up therefore rises from parity at $1^3$, where
setting up the job costs as much as the single MAC, to 457.2$\times$ at
$64^3$. In simulation, non-square and odd shapes behave the same way, for
example $5\times3\times7$: 20.9$\times$, $13\times16\times9$: 124.7$\times$,
$32\times1\times64$: 52.2$\times$ and $32\times64\times32$: 407.9$\times$.

\begin{table*}[t]
\caption{On-board cycle counts at 50\,MHz (square jobs $n\times n\times n$, all PASS)}\label{tab:perf}
\centering\footnotesize
\begin{tabular}{rrrrrrrr}
\toprule
$n$ & MACs & Software cycles & Accelerator cycles & Fetch & Compute + copy & Driver overhead & Speed-up \\
\midrule
1 & 1 & 118 & 118 & 14 & 20 & 84 & 1.0$\times$ \\
2 & 8 & 513 & 118 & 16 & 23 & 79 & 4.3$\times$ \\
4 & 64 & 3,127 & 135 & 20 & 35 & 80 & 23.2$\times$ \\
8 & 512 & 22,203 & 237 & 44 & 111 & 82 & 93.7$\times$ \\
16 & 4,096 & 168,163 & 747 & 140 & 527 & 80 & 225.1$\times$ \\
32 & 32,768 & 1,310,643 & 3,705 & 524 & 3,087 & 94 & 353.8$\times$ \\
64 & 262,144 & 10,352,467 & 22,643 & 2,060 & 20,495 & 88 & \textbf{457.2$\times$} \\
\bottomrule
\end{tabular}
\end{table*}

\begin{figure}[t]
\centering
\setlength{\unitlength}{1mm}
\begin{picture}(84,52)(-8,-6)
\put(0,0){\line(0,1){44}}\put(0,0){\line(1,0){76}}
\multiput(0,7)(0,7){6}{\line(-1,0){1}}
\put(-7,6.3){\tiny $10^2$}\put(-7,13.3){\tiny $10^3$}\put(-7,20.3){\tiny $10^4$}
\put(-7,27.3){\tiny $10^5$}\put(-7,34.3){\tiny $10^6$}\put(-7,41.3){\tiny $10^7$}
\put(3,0){\rule{3.5mm}{7.50mm}}\put(6.5,0){\color{gray}\rule{3.5mm}{7.50mm}}
\put(14,0){\rule{3.5mm}{11.97mm}}\put(17.5,0){\color{gray}\rule{3.5mm}{7.50mm}}
\put(25,0){\rule{3.5mm}{17.47mm}}\put(28.5,0){\color{gray}\rule{3.5mm}{7.91mm}}
\put(36,0){\rule{3.5mm}{23.42mm}}\put(39.5,0){\color{gray}\rule{3.5mm}{9.62mm}}
\put(47,0){\rule{3.5mm}{29.58mm}}\put(50.5,0){\color{gray}\rule{3.5mm}{13.11mm}}
\put(58,0){\rule{3.5mm}{35.82mm}}\put(61.5,0){\color{gray}\rule{3.5mm}{17.98mm}}
\put(69,0){\rule{3.5mm}{42.11mm}}\put(72.5,0){\color{gray}\rule{3.5mm}{23.49mm}}
\put(5,-3){\tiny $1^3$}\put(16,-3){\tiny $2^3$}\put(27,-3){\tiny $4^3$}
\put(38,-3){\tiny $8^3$}\put(48.5,-3){\tiny $16^3$}\put(59.5,-3){\tiny $32^3$}
\put(70.5,-3){\tiny $64^3$}
\put(3,41){\rule{2mm}{2mm}}\put(6,41.3){\tiny software (CPU)}
\put(3,37){\color{gray}\rule{2mm}{2mm}}\put(6,37.3){\tiny accelerator (end to end)}
\put(-8,46){\tiny cycles}\put(30,-6){\tiny matrix size $n\times n\times n$}
\end{picture}
\caption{Clock cycles on the board, log scale.}
\label{fig:cycles}
\end{figure}

\begin{figure}[t]
\centering
\setlength{\unitlength}{1mm}
\begin{picture}(84,52)(-8,-6)
\put(0,0){\line(0,1){46}}\put(0,0){\line(1,0){76}}
\multiput(0,9)(0,9){5}{\line(-1,0){1}}
\put(-6,8.3){\tiny 100}\put(-6,17.3){\tiny 200}\put(-6,26.3){\tiny 300}
\put(-6,35.3){\tiny 400}\put(-6,44.3){\tiny 500}
\put(4,0){\rule{5mm}{0.09mm}}\put(4.5,1.2){\tiny 1.0}
\put(15,0){\rule{5mm}{0.39mm}}\put(15.5,1.5){\tiny 4.3}
\put(26,0){\rule{5mm}{2.09mm}}\put(25.8,3.2){\tiny 23.2}
\put(37,0){\rule{5mm}{8.43mm}}\put(36.8,9.5){\tiny 93.7}
\put(48,0){\rule{5mm}{20.26mm}}\put(47.3,21.4){\tiny 225.1}
\put(59,0){\rule{5mm}{31.84mm}}\put(58.3,33.0){\tiny 353.8}
\put(70,0){\rule{5mm}{41.15mm}}\put(69.3,42.3){\tiny 457.2}
\put(5,-3){\tiny $1^3$}\put(16,-3){\tiny $2^3$}\put(27,-3){\tiny $4^3$}
\put(38,-3){\tiny $8^3$}\put(48.5,-3){\tiny $16^3$}\put(59.5,-3){\tiny $32^3$}
\put(70.5,-3){\tiny $64^3$}
\put(-8,48){\tiny speed-up ($\times$)}\put(30,-6){\tiny matrix size $n\times n\times n$}
\end{picture}
\caption{Speed-up of the accelerator over software on the same PicoRV32.}
\label{fig:speedup}
\end{figure}

\subsection{Analysis}\label{sec:analysis}

\emph{Where the CPU's cycles go.} With \texttt{-O2}, the inner loop of the
software kernel is seven instructions per MAC:
\begin{lstlisting}
loop: lb a4,0(a5)    # A[i][t]
      lb a1,0(a3)    # B[t][j]
      addi a5,a5,1   # next A element
      add a3,a3,a6   # next B element (stride N)
      mul a4,a4,a1
      add a2,a2,a4   # sum += product
      bne a7,a5,loop
\end{lstlisting}
Table~\ref{tab:cpu} gives the measured instruction counts. At $64^3$, the CPU
retires 7.11 instructions per MAC at a CPI of 5.55, which is 39.5 cycles per
MAC. PicoRV32 is not pipelined: every instruction passes through several
states, and every memory access, including each instruction fetch, waits one
clock for the block RAM. Only one of the seven instructions is the multiply.
The other six handle addressing, loading and looping, which the accelerator
replaces with counters and wiring.

\begin{table}[t]
\caption{Software kernel on PicoRV32 (measured on board)}\label{tab:cpu}
\centering\footnotesize
\begin{tabular}{rrrrr}
\toprule
$n$ & Instructions & CPI & Cycles/MAC & Instr./MAC \\
\midrule
1 & 29 & 4.07 & 118.0 & 29.0 \\
4 & 599 & 5.22 & 48.9 & 9.4 \\
16 & 30,599 & 5.50 & 41.1 & 7.5 \\
64 & 1,864,199 & 5.55 & 39.5 & 7.1 \\
\bottomrule
\end{tabular}
\end{table}

\emph{Why the speed-up is large.} The speed-up factors exactly into two
measured terms:
\begin{equation}
\frac{10{,}352{,}467}{22{,}643}
=\underbrace{\frac{\text{CPU cycles}}{\text{MAC}}}_{39.49}
\times\underbrace{\frac{\text{MACs}}{\text{accel.\ cycle}}}_{11.58}
=457.2 .
\end{equation}
The array's peak is 16 MACs per clock, so the bound is
$16\times39.5\approx632\times$. The accelerator reaches 72\% of it; the rest of
the time goes to fetch (2,060 cycles, i.e.\ the 2,048 words of $A$ and $B$),
copy-back ($\approx$4,100) and the driver (88). Part of the gain comes from the
weak baseline: a pipelined CPU at CPI 1 would still need about 7 cycles per MAC,
and the same accelerator would then give about 81$\times$.

\emph{Why the scratchpads matter.} Each beat needs one byte from four
different rows of $A$, which are four separate words in main RAM, plus four
bytes of $B$. That is at least five reads per clock on a single port. The banked
scratchpads deliver a beat in one read each. During compute, $A$ is also read
$N/4$ times and $B$ $M/4$ times, 16 times each at $64^3$, but each is fetched
from main RAM only once. Reading main RAM directly would take at least five
clocks per beat, or about 82,000 cycles instead of 20,500 at $64^3$, and the
speed-up would fall to roughly 120$\times$.

\emph{Real time and energy.} If each build runs at its own $f_{max}$, the
$64^3$ speed-up is still $457\times92/111.5\approx377\times$. At 50\,MHz, the
CPU-only system takes 207\,ms and about 47\,mJ, while the accelerated system
takes 0.45\,ms and about 0.11\,mJ: roughly 430$\times$ less energy, based on
Vivado's power estimates.

% ================================================================ VI
\section{Conclusions}\label{sec:concl}

We implemented a complete INT8 matrix-multiplication accelerator, from Booth
and carry-save adder cells to a memory-mapped peripheral of a RISC-V
processor, and ran it on a PYNQ-Z2. It occupies 8\% of the device's LUTs and no
DSPs, meets timing at 50\,MHz with an estimated $f_{max}$ of 92\,MHz, and computes a
$64\times64\times64$ product in 457$\times$ fewer cycles than the same processor
in software, verified element by element on the board.

Future work follows from our measurements. First, overlapping the phases:
storing finished rows of $C$ during compute would reduce a $16^3$ job from 667
to about 473 cycles, and a three-stage pipeline across jobs with
double-buffered scratchpads to about 399. Second, replicating the drain and
reset nets per row should recover most of the CPU-only clock rate. Third,
using the existing interrupt, a scaling study with $8\times8$ arrays, and PCPI
custom instructions.

% ================================================================ refs
\begin{thebibliography}{10}\footnotesize\setlength\itemsep{0pt}
\bibitem{kung} H.-T. Kung, ``Why systolic architectures?'' \emph{Computer},
vol.~15, no.~1, pp.~37--46, 1982.
\bibitem{tpu} N.~P. Jouppi \emph{et al.}, ``In-datacenter performance
analysis of a tensor processing unit,'' in \emph{Proc. 44th ISCA}, 2017,
pp.~1--12.
\bibitem{eyeriss} Y.-H. Chen \emph{et al.}, ``Eyeriss: An energy-efficient
reconfigurable accelerator for deep convolutional neural networks,''
\emph{IEEE J. Solid-State Circuits}, vol.~52, no.~1, pp.~127--138, 2016.
\bibitem{puscasu} A.~Pu\c{s}ca\c{s}u, C.~B. Ciobanu and O.~Buiu, ``Systolic
array matrix multiplication accelerator,'' in \emph{Proc. Int. Semiconductor
Conf. (CAS)}, 2024, pp.~207--210.
\bibitem{booth} A.~D. Booth, ``A signed binary multiplication technique,''
\emph{Quart. J. Mech. Appl. Math.}, vol.~4, no.~2, pp.~236--240, 1951.
\bibitem{macsorley} O.~L. MacSorley, ``High-speed arithmetic in binary
computers,'' \emph{Proc. IRE}, vol.~49, no.~1, pp.~67--91, 1961.
\bibitem{wallace} C.~S. Wallace, ``A suggestion for a fast multiplier,''
\emph{IEEE Trans. Electron. Comput.}, vol.~EC-13, no.~1, pp.~14--17, 1964.
\bibitem{inayat} K.~Inayat \emph{et al.}, ``Factored systolic arrays based on
radix-8 multiplication for machine learning acceleration,'' \emph{IEEE Trans.
VLSI Syst.}, 2024.
\bibitem{picorv32} C.~Wolf, ``PicoRV32: A size-optimized RISC-V CPU.''
[Online]. Available: https://github.com/YosysHQ/picorv32
\bibitem{hp} J.~L. Hennessy and D.~A. Patterson, \emph{Computer Architecture:
A Quantitative Approach}, 6th~ed. Morgan Kaufmann, 2019.
\end{thebibliography}

% ================================================================ appendices
\appendix
\renewcommand\thesection{\Alph{section}}
\section{Presentation Feedback}

Each question raised during the presentation is listed with our answer and
whether it was incorporated.

\emph{Q1. Is the RAM the FPGA's block RAM?} Yes. Every memory is on-chip
block RAM, and the board's DDR is not used: the main RAM uses 16 RAMB36 tiles,
and scratchpads A (4) and B (1) and the C store (4) use the accelerator's 9
(Section~\ref{sec:setup}). \textbf{Answered; stated in the report.}

\emph{Q2. Is the scratchpad useful when main RAM also has one cycle of
latency?} The latency is the same. The benefit is bandwidth (one read per beat
instead of at least five on one port), $16\times$ reuse of $A$ and $B$ at
$64^3$, and alignment and padding done once. Without the scratchpads, the
$64^3$ speed-up would fall to roughly 120$\times$
(Section~\ref{sec:analysis}). \textbf{Answered; quantitative argument added;
design unchanged.}

\emph{Q3. Why such a huge speed-up with so little hardware?} The speed-up is
CPU cycles per MAC (39.5) times accelerator MACs per cycle (11.6). The
non-pipelined CPU needs seven instructions at CPI 5.5 per MAC. Against a CPI-1
processor, the gain would be about 81$\times$. The accelerator is 3.6$\times$
the CPU's LUTs, but only 8\% of the chip. \textbf{Incorporated: analysis with
measured CPI (Section~\ref{sec:analysis}).}

\emph{Q4. Explain one line of the speed-up screenshot.} The line was
\texttt{bench 16x16x16: software 168778 cycles, accelerator 842 cycles (FETCH
140 + COMPUTE 527), speed-up 200.4x}. \emph{Software} is \texttt{rdcycle}
around the triple loop, and \emph{accelerator} is \texttt{rdcycle} around the
whole driver call. FETCH is $(256+256)$ bytes, or 128 words, plus setup.
COMPUTE is 268 array cycles (16 tiles of 16 beats, plus pipeline fill) plus
259 copy-back cycles. The remaining 175 cycles are register writes, START and
polling. \textbf{Answered; later firmware reports copy-back separately.}

\emph{Q5. Do you know which part contributes to the CPU time?} At the
presentation, we had only total cycles. We then added \texttt{rdinstret} to
the firmware and analysed the disassembly: 7 instructions per MAC, of which two
are loads and one is the multiply, at CPI 5.55 (Table~\ref{tab:cpu}).
\textbf{Incorporated: new measurement in the firmware and scripts.}

\emph{Q6. Have you checked how many cycles the normal CPU takes?} Yes. Every
speed-up uses the software loop on the same CPU. We also ran a CPU-only
bitstream through the same sweep, and it gives identical software cycle
counts. \textbf{Incorporated: CPU-only build and sweep.}

\emph{Q7. Have you done synthesis and a utilization report of the CPU
alone?} Not at the presentation. Since then, we have built and reported the
CPU-only system: 2,776 LUTs, 20 BRAM and 4 DSPs for the whole chip, with
PicoRV32 itself using 1,217 LUTs. It reaches $f_{max}\approx111.5$\,MHz at
0.228\,W (Tables~\ref{tab:util}--\ref{tab:blocks}). \textbf{Incorporated.}

\section{Disclosure of LLM Assistance}

We used an LLM (Claude) during the project; overall usage was moderate.
Table~\ref{tab:llm} lists what it was used for.

\begin{table}[h]
\caption{Use of LLM assistance}\label{tab:llm}
\centering\footnotesize
\begin{tabularx}{\columnwidth}{lLl}
\toprule
Area & How & Amount \\
\midrule
RTL coding & We specified behaviour and ports; the LLM wrote Verilog, and we reviewed each piece & Low \\
Verification & Testbenches, golden model, vectors and regression scripts from edge cases we specified & Heavy \\
Design decisions & Options looked up; decisions and reasons logged by us & Low \\
Concepts & Systolic dataflow, hazards, timing reports & Moderate \\
FPGA tooling & Vivado scripts, reading reports & Heavy \\
Visualizers & HTML cycle-by-cycle visualizers, diagrams & Heavy \\
Reports, slides & Drafted from our inputs, including this report; we checked content and numbers & Heavy \\
\bottomrule
\end{tabularx}
\end{table}

\section{Contributions}

\emph{Aditi Dure:} file hierarchy and reference survey, clocked Booth/CSA PE,
tile sequencer and four-bank C store, design-decision log, Vivado synthesis
and timing fixes, PYNQ-Z2 bring-up, interactive board flow, CPU-only baseline
and CPI measurement.
\emph{Krishna Kaustubh:} repository setup, $4\times4$ array, visualizers,
dual-port RAM, scratchpads, fetch unit, C copier, registers, job controller,
\texttt{soc\_top} integration and firmware, PYNQ-Z2 \texttt{fpga/} flow.

\section{Reproducing the Results}

The code is organised as follows: \texttt{rtl/} (Verilog), \texttt{tb/}
(testbenches), \texttt{golden/} (Python model), \texttt{fw/} (firmware),
\texttt{fpga/} (board top, constraints, Vivado scripts), \texttt{scripts/},
\texttt{visualizer/} and \texttt{docs/}. Running it requires Icarus Verilog,
Python~3 and Make; the RISC-V GCC toolchain and Vivado are needed only for
firmware and FPGA builds.
\begin{lstlisting}
make check              # 40 testbenches
make fw && make test-soc_top
make test-soc-full      # large jobs
make fpga-run-sim       # board flow in Icarus
make fpga-baseline-sim  # CPU-only build
vivado -mode batch -source fpga/build.tcl
vivado -mode batch -source fpga/build.tcl -tclargs baseline
# Vivado Tcl shell, board connected:
source fpga/run.tcl
run_program             # or: run_program baseline
sweep {1 2 4 8 16 32 64}
\end{lstlisting}

\end{document}
